\documentclass{amsart}
% For dealing with references we use the comment environment
\usepackage{verbatim}
\newenvironment{reference}{\comment}{\endcomment}
%\newenvironment{reference}{}{}
\newenvironment{slogan}{\comment}{\endcomment}
\newenvironment{history}{\comment}{\endcomment}

% For commutative diagrams we use Xy-pic
\usepackage[all]{xy}

% We use 2cell for 2-commutative diagrams.
\xyoption{2cell}
\UseAllTwocells

% We use multicol for the list of chapters between chapters
\usepackage{multicol}

% This is generally recommended for better output
\usepackage{lmodern}
\usepackage[T1]{fontenc}

% For cross-file-references

% Package for hypertext links:
\usepackage{hyperref}

% For any local file, say "hello.tex" you want to link to please

% Theorem environments.
%
\theoremstyle{plain}
\newtheorem{theorem}[subsection]{Theorem}
\newtheorem{proposition}[subsection]{Proposition}
\newtheorem{lemma}[subsection]{Lemma}

\theoremstyle{definition}
\newtheorem{definition}[subsection]{Definition}
\newtheorem{example}[subsection]{Example}
\newtheorem{exercise}[subsection]{Exercise}
\newtheorem{situation}[subsection]{Situation}

\theoremstyle{remark}
\newtheorem{remark}[subsection]{Remark}
\newtheorem{remarks}[subsection]{Remarks}

\numberwithin{equation}{subsection}

% Macros
%
\def\lim{\mathop{\mathrm{lim}}\nolimits}
\def\colim{\mathop{\mathrm{colim}}\nolimits}
\def\Spec{\mathop{\mathrm{Spec}}}
\def\Hom{\mathop{\mathrm{Hom}}\nolimits}
\def\Ext{\mathop{\mathrm{Ext}}\nolimits}
\def\SheafHom{\mathop{\mathcal{H}\!\mathit{om}}\nolimits}
\def\SheafExt{\mathop{\mathcal{E}\!\mathit{xt}}\nolimits}
\def\Sch{\mathit{Sch}}
\def\Mor{\mathop{\mathrm{Mor}}\nolimits}
\def\Ob{\mathop{\mathrm{Ob}}\nolimits}
\def\Sh{\mathop{\mathit{Sh}}\nolimits}
\def\NL{\mathop{N\!L}\nolimits}
\def\CH{\mathop{\mathrm{CH}}\nolimits}
\def\proetale{{pro\text{-}\acute{e}tale}}
\def\etale{{\acute{e}tale}}
\def\QCoh{\mathit{QCoh}}
\def\Ker{\mathop{\mathrm{Ker}}}
\def\Im{\mathop{\mathrm{Im}}}
\def\Coker{\mathop{\mathrm{Coker}}}
\def\Coim{\mathop{\mathrm{Coim}}}

% Boxtimes
%
\DeclareMathSymbol{\boxtimes}{\mathbin}{AMSa}{"02}

%
% Macros for moduli stacks/spaces
%
\def\QCohstack{\mathcal{QC}\!\mathit{oh}}
\def\Cohstack{\mathcal{C}\!\mathit{oh}}
\def\Spacesstack{\mathcal{S}\!\mathit{paces}}
\def\Quotfunctor{\mathrm{Quot}}
\def\Hilbfunctor{\mathrm{Hilb}}
\def\Curvesstack{\mathcal{C}\!\mathit{urves}}
\def\Polarizedstack{\mathcal{P}\!\mathit{olarized}}
\def\Complexesstack{\mathcal{C}\!\mathit{omplexes}}
% \Pic is the operator that assigns to X its picard group, usage \Pic(X)
% \Picardstack_{X/B} denotes the Picard stack of X over B
% \Picardfunctor_{X/B} denotes the Picard functor of X over B
\def\Pic{\mathop{\mathrm{Pic}}\nolimits}
\def\Picardstack{\mathcal{P}\!\mathit{ic}}
\def\Picardfunctor{\mathrm{Pic}}
\def\Deformationcategory{\mathcal{D}\!\mathit{ef}}

\setlength{\emergencystretch}{3em}
\begin{document}
\title[Stacks Project, Tag 05A9]{1509 --- Projectivity of arbitrary modules descends along any faithfully flat ring extension}
\maketitle
\noindent Copyright (C) 2005--2025 Johan de Jong. Stacks Project authors. Source: \href{https://stacks.math.columbia.edu/tag/05A9}{Tag 05A9}.

\noindent Editorial difficulty note (not author text). Difficulty H2: The proof descends countable projective summands through tensor-kernel stabilization, builds adapted countable submodules compatible with the given base-changed summands, and assembles them by a transfinite split filtration. Descent of unrestricted infinitely generated projectivity requires substantially more than ordinary finite-presentation descent..

\noindent Editorial provenance note (not author text). History: excerpt from revision \texttt{a0444\allowbreak 6e57e\allowbreak c1fbc\allowbreak 252a8\allowbreak 71afc\allowbreak ec775\allowbreak 2fb28\allowbreak 07b14}, algebra.tex, lines 22461--22515. Reference presentation was adapted; original-author context and core support blocks were retained or added. The mathematical wording is unchanged. Licensed under GNU FDL 1.2 or later; no invariant sections or cover texts. See ../licenses/COPYING. Context blocks reproduce author definitions and supporting results; they are not additional counted problems.

\begin{definition}
\label{definition-devissage}
Let $M$ be an $R$-module.  A {\it direct sum d\'evissage} of $M$ is a family
of submodules $(M_{\alpha})_{\alpha \in S}$, indexed by an ordinal $S$ and
increasing (with respect to inclusion), such that:
\begin{enumerate}
\item[(0)] $M_0 = 0$;
\item[(1)] $M = \bigcup_{\alpha} M_{\alpha}$;
\item[(2)] if $\alpha \in S$ is a limit ordinal, then $M_{\alpha} =
\bigcup_{\beta < \alpha} M_{\beta}$;
\item[(3)] if $\alpha + 1 \in S$, then $M_{\alpha}$ is a direct summand of
$M_{\alpha + 1}$.
\end{enumerate}
If moreover
\begin{enumerate}
\item[(4)] $M_{\alpha + 1}/M_{\alpha}$ is countably generated for
$\alpha + 1 \in S$,
\end{enumerate}
then $(M_{\alpha})_{\alpha \in S}$ is called a {\it Kaplansky d\'evissage}
of $M$.
\end{definition}

\begin{definition}
\label{definition-mittag-leffler-module}
Let $M$ be an $R$-module.  We say that $M$ is {\it Mittag-Leffler} if the
equivalent conditions of
Proposition \ref{proposition-ML-characterization}
hold.
\end{definition}

\begin{proposition}
\label{proposition-ML-characterization}
Let $M$ be an $R$-module.  Let $(M_i, f_{ij})$ be a directed system of finitely
presented $R$-modules, indexed by $I$, such that $M = \colim M_i$.  Let
$f_i:
M_i \to M$ be the canonical map.  The following are equivalent:
\begin{enumerate}
\item For every finitely presented $R$-module $P$ and module map $f: P
\to M$, there exists a finitely presented $R$-module $Q$ and a module
map $g: P \to Q$ such that $g$ and $f$ dominate each other, i.e.,
$\Ker(f \otimes_R \text{id}_N) = \Ker(g \otimes_R \text{id}_N)$
for every $R$-module $N$.
\item For each $i \in I$, there exists $j \geq i$ such that $f_{ij}: M_i
\to M_j$ dominates $f_i: M_i \to M$.
\item For each $i \in I$, there exists $j \geq i$ such that $f_{ij}: M_i
\to M_j$ factors through $f_{ik}: M_i \to M_k$ for all $k \geq
i$.
\item For every $R$-module $N$, the inverse system
$(\Hom_R(M_i, N), \Hom_R(f_{ij}, N))$ is Mittag-Leffler.
\item For $N = \prod_{s \in I} M_s$, the inverse system
$(\Hom_R(M_i, N), \Hom_R(f_{ij}, N))$ is Mittag-Leffler.
\end{enumerate}
\end{proposition}

\begin{definition}
\label{definition-domination}
Let $f: M \to N$ and $g: M \to M'$ be maps of $R$-modules.
Then we say $g$ {\it dominates} $f$ if for any $R$-module $Q$, we have $\Ker(f
\otimes_R \text{id}_Q) \subset \Ker(g \otimes_R \text{id}_Q)$.
\end{definition}

\begin{definition}
\label{definition-universally-injective}
Let $f: M \to N$ be a map of $R$-modules.  Then $f$ is called
{\it universally injective} if for every $R$-module $Q$, the map $f
\otimes_R \text{id}_Q: M \otimes_R Q \to N \otimes_R Q$
is injective.  A sequence $0 \to M_1 \to M_2 \to M_3
\to 0$ of $R$-modules is called {\it universally exact} if it is exact
and $M_1 \to M_2$ is universally injective.
\end{definition}

\begin{lemma}
\label{lemma-adapted-submodule}
Let $R \to S$ be a ring map, and let $M$ be an $R$-module.  Suppose $M
\otimes_R S = \bigoplus_{i \in I} Q_i$ is a direct sum of countably generated
$S$-modules $Q_i$.  If $N$ is a countably generated submodule of $M$, then
there is a countably generated submodule $N'$ of $M$ such that $N' \supset N$
and $\Im(N' \otimes_R S \to M \otimes_R S) =
\bigoplus_{i \in I'} Q_i$ for some subset $I' \subset I$.
\end{lemma}

\begin{proof}
Let $N'_0 = N$.  We construct by induction an increasing sequence of countably
generated submodules $N'_{\ell} \subset M$ for $\ell = 0, 1, 2, \ldots$
such that: if $I'_{\ell}$ is the set of $i \in I$ such that the projection of
$\Im(N'_{\ell} \otimes_R S \to M \otimes_R S)$ onto $Q_i$ is
nonzero, then $\Im(N'_{\ell + 1} \otimes_R S \to M \otimes_R
S)$ contains $Q_i$ for all $i \in I'_{\ell}$.  To construct $N'_{\ell + 1}$
from $N'_\ell$, let $Q$ be the sum of (the countably many) $Q_i$ for
$i \in I'_{\ell}$, choose $P$ as in Lemma
\ref{lemma-lift-countably-generated-submodule}, and then let $N'_{\ell + 1} =
N'_{\ell} + P$.  Having constructed the $N'_{\ell}$, just take $N' =
\bigcup_{\ell} N'_{\ell}$ and $I' = \bigcup_{\ell} I'_{\ell}$.
\end{proof}


\begin{lemma}
\label{lemma-lift-countably-generated-submodule}
Let $R \to S$ be a ring map, let $M$ be an $R$-module, and let $Q$ be a
countably generated $S$-submodule of $M \otimes_R S$.  Then there exists a
countably generated $R$-submodule $P$ of $M$ such that
$\Im(P \otimes_R S \to M \otimes_R S)$ contains $Q$.
\end{lemma}

\begin{proof}
Let $y_1, y_2, \ldots$ be generators for $Q$ and write $y_j = \sum_k x_{jk}
\otimes s_{jk}$ for some $x_{jk} \in M$ and $s_{jk} \in S$. Then take $P$ be
the submodule of $M$ generated by the $x_{jk}$.
\end{proof}


\begin{lemma}
\label{lemma-ffdescent-ML}
\begin{reference}
Email from Juan Pablo Acosta Lopez dated 12/20/14.
\end{reference}
Let $R \to S$ be a faithfully flat ring map. Let $M$ be an $R$-module. If the
$S$-module $M \otimes_R S$ is Mittag-Leffler, then $M$ is Mittag-Leffler.
\end{lemma}

\begin{proof}
Write $M = \colim_{i\in I} M_i$ as a directed colimit
of finitely presented $R$-modules $M_i$.
Using Proposition \ref{proposition-ML-characterization}, we see that we
have to prove that for each $i \in I$ there exists $i \leq j$, $j\in I$
such that $M_i\rightarrow M_j$ dominates $M_i\rightarrow M$.

\medskip\noindent
Take $N$ the pushout
$$
\xymatrix{
M_i  \ar[r] \ar[d] & M_j \ar[d] \\
M \ar[r] & N
}
$$
Then the lemma is equivalent to the existence of $j$ such that
$M_j\rightarrow N$ is universally injective, see
Lemma \href{https://stacks.math.columbia.edu/tag/0AUM}{lemma domination universally injective (Tag 0AUM)}.
Observe that the tensorization by $S$
$$
\xymatrix{
M_i\otimes_R S  \ar[r] \ar[d] & M_j\otimes_R S \ar[d] \\
M\otimes_R S \ar[r] & N\otimes_R S
}
$$
Is a pushout diagram.  So because
$M \otimes_R S = \colim_{i\in I} M_i \otimes_R S$
expresses $M\otimes_R S$ as a colimit of $S$-modules of
finite presentation, and $M\otimes_R S$ is Mittag-Leffler,
there exists $j \geq i$ such that $M_j\otimes_R S\rightarrow N\otimes_R S$
is universally injective. So using that $R\rightarrow S$ is faithfully flat
we conclude that $M_j\rightarrow N$ is universally injective too.
\end{proof}


\begin{lemma}
\label{lemma-ffdescent-countable}
Let $R \to S$ be a faithfully flat ring map. Let $M$ be an $R$-module. If
the $S$-module $M \otimes_R S$ is countably generated, then $M$
is countably generated.
\end{lemma}

\begin{proof}
Say $M \otimes_R S$ is generated by the elements
$y_i$, $i = 1, 2, 3, \ldots$. Write $y_i = \sum_{j = 1, \ldots, n_i}
x_{ij} \otimes s_{ij}$ for some $n_i \geq 0$, $x_{ij} \in M$
and $s_{ij} \in S$. Denote $M' \subset M$ the submodule generated
by the countable collection of elements $x_{ij}$. Then
$M' \otimes_R S \to M \otimes_R S$ is surjective as the
image contains the generators $y_i$. Since $S$ is faithfully flat
over $R$ we conclude that $M' = M$ as desired.
\end{proof}


\begin{lemma}
\label{lemma-ffdescent-countable-projectivity}
Let $R \to S$ be a faithfully flat ring map.  Let $M$ be an $R$-module.
 If the $S$-module $M \otimes_R S$ is countably generated and projective,
then $M$ is countably generated and projective.
\end{lemma}

\begin{proof}
Follows from Lemmas \href{https://stacks.math.columbia.edu/tag/03C4}{lemma descend properties modules (Tag 03C4)},
\ref{lemma-ffdescent-ML}, and \ref{lemma-ffdescent-countable} and
Theorem \href{https://stacks.math.columbia.edu/tag/059Z}{theorem projectivity characterization (Tag 059Z)}.
\end{proof}


\begin{theorem}
\label{theorem-kaplansky-direct-sum}
Suppose $M$ is a direct sum of countably generated $R$-modules.  If $P$ is a
direct summand of $M$, then $P$ is also a direct sum of countably generated
$R$-modules.
\end{theorem}

\begin{proof}
Write $M = P \oplus Q$.  We are going to construct a Kaplansky d\'evissage
$(M_{\alpha})_{\alpha \in S}$ of $M$ which, in addition to the defining
properties (0)-(4), satisfies:
\begin{enumerate}
\item[(5)] Each $M_{\alpha}$ is a direct summand of $M$;
\item[(6)] $M_{\alpha} = P_{\alpha} \oplus Q_{\alpha}$, where $P_{\alpha} =P
\cap M_{\alpha}$ and $Q_\alpha = Q \cap M_{\alpha}$.
\end{enumerate}
(Note: if properties (0)-(2) hold, then in fact property (3) is equivalent to
property (5).)

\medskip\noindent
To see how this implies the theorem, it is enough to show that
$(P_{\alpha})_{\alpha \in S}$ forms a Kaplansky d\'evissage of $P$.  Properties
(0), (1), and (2) are clear.  By (5) and (6) for $(M_{\alpha})$, each
$P_{\alpha}$ is a direct summand of $M$.  Since $P_{\alpha} \subset P_{\alpha +
1}$, this implies $P_{\alpha}$ is a direct summand of $P_{\alpha + 1}$; hence
(3) holds for $(P_{\alpha})$.  For (4), note that
$$
M_{\alpha + 1}/M_{\alpha} \cong P_{\alpha + 1}/P_{\alpha} \oplus
Q_{\alpha + 1}/Q_{\alpha},
$$
so $P_{\alpha + 1}/P_{\alpha}$ is countably generated because this is true of
$M_{\alpha + 1}/M_{\alpha}$.

\medskip\noindent
It remains to construct the $M_{\alpha}$.  Write $M = \bigoplus_{i \in I} N_i$
where each $N_i$ is a countably generated $R$-module.  Choose a well-ordering
of $I$.  By transfinite recursion we are going to define an increasing family
of submodules $M_{\alpha}$ of $M$, one for each ordinal $\alpha$, such that
$M_{\alpha}$ is a direct sum of some subset of the $N_i$.

\medskip\noindent
For $\alpha = 0$ let $M_{0} = 0$.  If $\alpha$ is a limit ordinal and
$M_{\beta}$ has been defined for all $\beta < \alpha$, then define $M_{\alpha}
= \bigcup_{\beta < \alpha} M_{\beta}$.  Since each $M_{\beta}$ for $\beta <
\alpha$ is a direct sum of a subset of the $N_i$, the same will be true of
$M_{\alpha}$.  If $\alpha + 1$ is a successor ordinal and $M_{\alpha}$ has been
defined, then define $M_{\alpha + 1}$ as follows.  If $M_{\alpha} = M$, then let
$M_{\alpha + 1} = M$.  If not, choose the smallest $j \in I$ such that $N_j$ is
not contained in $M_{\alpha}$.  We will construct an infinite matrix $(x_{mn}),
m, n = 1, 2, 3, \ldots$ such that:
\begin{enumerate}
\item $N_j$ is contained in the submodule of $M$ generated by the entries
$x_{mn}$;
\item if we write any entry $x_{k\ell}$ in terms of its $P$- and
$Q$-components, $x_{k\ell} = y_{k\ell} + z_{k\ell}$, then the matrix $(x_{mn})$
contains a set of generators for each $N_i$ for which $y_{k\ell}$ or
$z_{k\ell}$ has nonzero component.
\end{enumerate}
Then we define $M_{\alpha + 1}$ to be the submodule of $M$ generated by
$M_{\alpha}$ and all $x_{mn}$; by property (2) of the matrix $(x_{mn})$,
$M_{\alpha + 1}$ will be a direct sum of some subset of the $N_i$.
To construct the matrix $(x_{mn})$, let $x_{11}, x_{12}, x_{13}, \ldots$
be a countable set of generators for $N_j$. Then if
$x_{11} = y_{11} + z_{11}$ is the decomposition into $P$- and
$Q$-components, let $x_{21}, x_{22}, x_{23}, \ldots$ be a countable
set of generators for the sum of the $N_i$ for which $y_{11}$ or $z_{11}$ have
nonzero component.  Repeat this process on $x_{12}$ to get elements $x_{31},
x_{32}, \ldots$, the third row of our matrix.  Repeat on $x_{21}$ to get the
fourth row, on $x_{13}$ to get the fifth, and so on, going down along
successive anti-diagonals as indicated below:
$$
\left(
\vcenter{
\xymatrix@R=2mm@C=2mm{
x_{11} & x_{12} \ar[dl] & x_{13} \ar[dl] & x_{14} \ar[dl] & \ldots  \\
x_{21} & x_{22} \ar[dl] & x_{23} \ar[dl] & \ldots  \\
x_{31} & x_{32} \ar[dl] & \ldots \\
x_{41} & \ldots \\
\ldots
}
}
\right).
$$

\medskip\noindent
Transfinite induction on $I$ (using the fact that we constructed
$M_{\alpha + 1}$
to contain $N_j$ for the smallest $j$ such that $N_j$ is not contained in
$M_{\alpha}$) shows that for each $i \in I$, $N_i$ is contained in some
$M_{\alpha}$.  Thus, there is some large enough ordinal $S$ satisfying: for
each $i \in I$ there is $\alpha \in S$ such that $N_i$ is contained in
$M_{\alpha}$.  This means $(M_{\alpha})_{\alpha \in S}$ satisfies property (1)
of a Kaplansky d\'evissage of $M$.  The family $(M_{\alpha})_{\alpha \in S}$
moreover satisfies the other defining properties, and also (5) and (6) above:
properties (0), (2), (4), and (6) are clear by construction;  property (5) is
true because each $M_{\alpha}$ is by construction a direct sum of some $N_i$;
and (3) is implied by (5) and the fact that $M_{\alpha} \subset M_{\alpha + 1}$.
\end{proof}


\begin{lemma}
\label{lemma-countgen-projective}
Let $M$ be an $R$-module.  If $M$ is flat, Mittag-Leffler, and countably
generated, then $M$ is projective.
\end{lemma}

\begin{proof}
By Lazard's theorem (Theorem \href{https://stacks.math.columbia.edu/tag/058G}{theorem lazard (Tag 058G)}), we can write $M =
\colim_{i
\in I} M_i$ for a directed system of finite free $R$-modules $(M_i, f_{ij})$
indexed by a set $I$.  By Lemma \href{https://stacks.math.columbia.edu/tag/059W}{lemma ML countable colimit (Tag 059W)}, we may assume
$I$ is countable.  Now let
$$
0 \to N_1 \to N_2 \to N_3 \to 0
$$
be an exact sequence of $R$-modules.  We must show that applying
$\Hom_R(M, -)$
preserves exactness.  Since $M_i$ is finite free,
$$
0 \to \Hom_R(M_i, N_1) \to \Hom_R(M_i, N_2) \to
\Hom_R(M_i, N_3) \to 0
$$
is exact for each $i$.  Since $M$ is Mittag-Leffler, $(\Hom_R(M_i,
N_{1}))$ is
a Mittag-Leffler inverse system.  So by Lemma \href{https://stacks.math.columbia.edu/tag/0598}{lemma ML exact sequence (Tag 0598)},
$$
0 \to \lim_{i \in I} \Hom_R(M_i, N_1) \to
\lim_{i \in I} \Hom_R(M_i, N_2) \to
\lim_{i \in I} \Hom_R(M_i, N_3) \to 0
$$
is exact.  But for any $R$-module $N$ there is a functorial isomorphism
$\Hom_R(M, N) \cong \lim_{i \in I} \Hom_R(M_i, N)$, so
$$
0 \to \Hom_R(M, N_1) \to \Hom_R(M, N_2) \to
\Hom_R(M, N_3) \to 0
$$
is exact.
\end{proof}


\begin{theorem}
\label{theorem-ffdescent-projectivity}
Let $R \to S$ be a faithfully flat ring map.  Let $M$ be an $R$-module.
 If the $S$-module $M \otimes_R S$ is projective, then $M$ is projective.
\end{theorem}

\begin{proof}
We are going to construct a Kaplansky d\'evissage of $M$ to show that it is a
direct sum of projective modules and hence projective.  By Theorem
\href{https://stacks.math.columbia.edu/tag/058Y}{theorem projective direct sum (Tag 058Y)} we can write $M \otimes_R S =
\bigoplus_{i \in I} Q_i$ as a direct sum of countably generated $S$-modules
$Q_i$.  Choose a well-ordering on $M$.  Using transfinite recursion we are going
to define an increasing family of submodules $M_{\alpha}$ of $M$, one for each
ordinal $\alpha$, such that $M_{\alpha} \otimes_R S$ is a direct sum of some
subset of the $Q_i$.

\medskip\noindent
For $\alpha = 0$ let $M_0 = 0$.  If $\alpha$ is a limit ordinal and $M_{\beta}$
has been defined for all $\beta < \alpha$, then define $M_\alpha =
\bigcup_{\beta < \alpha} M_{\beta}$.  Since each $M_{\beta} \otimes_R S$ for
$\beta < \alpha$ is a direct sum of a subset of the $Q_i$, the same will be
true of $M_{\alpha} \otimes_R S$.  If $\alpha + 1$ is a successor ordinal and
$M_{\alpha}$ has been defined, then define $M_{\alpha + 1}$ as follows.  If
$M_{\alpha} = M$, then let $M_{\alpha +1} = M$.  Otherwise choose the smallest
$x \in M$ (with respect to the fixed well-ordering) such that $x \notin
M_{\alpha}$. Since $S$ is flat over $R$, $(M/M_{\alpha}) \otimes_R S = M
\otimes_R S/M_{\alpha} \otimes_R S$, so since $M_{\alpha} \otimes_R S$ is
a direct sum of some $Q_i$, the same is true of $(M/M_{\alpha}) \otimes_R
S$.   By Lemma \ref{lemma-adapted-submodule}, we can find a countably generated
$R$-submodule $P$ of $M/M_{\alpha}$ containing the image of $x$ in
$M/M_{\alpha}$ and such that $P \otimes_R S$ (which equals $\Im(P
\otimes_R S \to M \otimes_R S)$ since $S$ is flat over $R$) is a
direct sum of some $Q_i$. Since $M \otimes_R S = \bigoplus_{i \in I} Q_i$
is projective and projectivity passes to direct summands, $P \otimes_R S$ is
also projective.  Thus by Lemma \ref{lemma-ffdescent-countable-projectivity},
$P$ is projective.  Finally we define $M_{\alpha + 1}$ to be the preimage of $P$
in $M$, so that $M_{\alpha + 1}/M_{\alpha} = P$ is countably generated and
projective.  In particular $M_{\alpha}$ is a direct summand of $M_{\alpha + 1}$
since projectivity of $M_{\alpha + 1}/M_{\alpha}$ implies the sequence $0
\to M_{\alpha} \to M_{\alpha + 1} \to
M_{\alpha + 1}/M_{\alpha} \to 0$ splits.

\medskip\noindent
Transfinite induction on $M$ (using the fact that we constructed
$M_{\alpha + 1}$ to contain the smallest $x \in M$ not contained in
$M_{\alpha}$) shows that
each $x \in M$ is contained in some $M_{\alpha}$.  Thus, there is some large
enough ordinal $S$ satisfying: for each $x \in M$ there is $\alpha \in S$ such
that $x \in M_{\alpha}$.  This means $(M_{\alpha})_{\alpha \in S}$ satisfies
property (1) of a Kaplansky d\'evissage of $M$.  The other properties are clear
by construction.  We conclude
$M = \bigoplus_{\alpha + 1 \in S} M_{\alpha + 1}/M_{\alpha}$.
Since each $M_{\alpha + 1}/M_{\alpha}$ is projective
by construction, $M$ is projective.
\end{proof}
\end{document}
